% id: 118-01
% book: 베이직쎈 확률과 통계 · 06 이항분포와 정규분포
% section: 유형 074 이항분포와 정규분포의 관계
% figure: none
% answer: 
% answer_source: 
% variant_level: 0 (원본 전사)
% 이 파일은 build.mjs 가 items.json 에서 만든다. 고칠 때는 items.json 을 고친다.
\smash{\raisebox{4.5mm}{\dmkichul{베이직쎈 확률과 통계 118쪽 01번}}}
확률변수 $X$가 이항분포 $\mathrm{B}\left(180{,}\mkern3mu\mathopen{}\dfrac{1}{6}\right)$을 따르면 $X$는 근사적으로 정규분포를 따른다. 이때 $X$가 따르는 정규분포를 바르게 나타낸 것은?
\begin{choices32}{$\mathrm{N}(30{,}\mkern3mu\mathopen{}3^2)$}{$\mathrm{N}(30{,}\mkern3mu\mathopen{}5^2)$}{$\mathrm{N}(30{,}\mkern3mu\mathopen{}9^2)$}{$\mathrm{N}(60{,}\mkern3mu\mathopen{}3^2)$}{$\mathrm{N}(60{,}\mkern3mu\mathopen{}5^2)$}\end{choices32}
